\chapter{Introduction}\label{ch:intro}

\micro{} is a small, fast miniKanren. Its terms are first-order Scheme data, its unifier is first-order
unification with an occurs check, and its goals are conjunctions, disjunctions and existentials. In
logical terms it runs first-order Horn clauses (\cref{sec:horn-micro}). λProlog, the language of Miller and
Nadathur's \emph{Programming with Higher-Order Logic}~\cite{MillerNadathur2012}, runs a much larger logic.
Goals may say \enquote{for every new name $x$} and \enquote{assuming $D$}, terms may contain λ-binders, and
unification works modulo α, β and η. With these three features, programs over syntax with binders (type
checkers, evaluators, provers, program transformations) contain no code for renaming, fresh names or
capture-avoiding substitution, because the logic does that work.

This document asks what it would take to give \micro{} the same power, and it ends with a concrete answer:
a module, \texttt{(aux kanren lambda)}, prototyped, tested and benchmarked. This chapter says what the
document is for (\cref{sec:intro-goal}), runs one example through the whole story
(\cref{sec:intro-brief}), maps the chapters (\cref{sec:intro-map}), proposes three ways to read them
(\cref{sec:intro-paths}), and fixes the conventions (\cref{sec:intro-conventions}).

% ==================================================================================================
\section{What this document is for}\label{sec:intro-goal}

\paragraph{The goal.} The goal is to bring three ideas of λProlog into \micro:
\begin{enumerate}[nosep]
  \item \emergencystretch=2em \textbf{hereditary Harrop goals}, that is the universal goal $\forall x\,G$ (\lp{pi x\ G}, proved
        with a new scoped constant) and the hypothetical goal $D\imp G$ (\lp{D => G}, proved with the
        program temporarily extended by the clause $D$);
  \item \textbf{λ-tree syntax}: object-level binders represented by meta-level λ-ab\-strac\-tions, so that
        α-equivalence is built in and object-level substitution is meta-level β-reduction;
  \item \textbf{higher-order pattern unification} ($\Llam$): unification of λ-terms restricted to the
        fragment where it is decidable and has most general unifiers, with the other equations delayed.
\end{enumerate}
\noindent The three are not independent. \lp{pi} is what lets a program walk under a λ-binder, and
pattern unification is what lets it build a λ-term back. The book's programs use all three together.

\paragraph{The sources.} There are two. The first is the book itself, which Part~I summarizes chapter by
chapter, with its core examples run on \micro{} wherever that is possible today. The second is the
miniKanren literature on these features, which Part~II surveys. Its center is
Ma and Friedman's \emph{A New Higher-order Unification Algorithm for λKanren}~\cite{MaFriedman2021}, read
together with its predecessor λKanren~\cite{MaLuFriedman2020}, nominal unification and
αKanren~\cite{UrbanPittsGabbay2004,ByrdFriedman2007}, the eigenvariable and negation
work~\cite{HemannByrdFriedman2013,Moiseenko2019,JinEtAl2021}, and the techniques of real λProlog systems:
universe levels~\cite{NadathurJayaramanKwon1995}, Miller's and Nipkow's pattern
unification~\cite{Miller1991,Nipkow1993}, delayed equations~\cite{DowekHardinKirchnerPfenning1996},
Teyjus~\cite{NadathurLinnell2005} and ELPI~\cite{DunchevEtAl2015}.

\paragraph{The outcome.} Part~III is a proposal in two pieces.
\begin{itemize}[nosep]
  \item \emergencystretch=2em A \textbf{generic extension protocol} in \micro{} (\texttt{micro.diff}, $+235/{-}41$ lines,
        comments included; 214 of the added lines are not blank). The state gains one slot that \micro{} never reads, and a registry of hooks opens the
        points that are closed today: the unifier's structural clause, the occurs check, the step after each
        \scm{=°}, the disequality, type and absento stores, and reification. Nothing in the edit mentions λ,
        levels or clauses.
  \item A \textbf{client module}, \texttt{(aux kanren lambda)}, 1255 lines, which registers one extension.
        It provides λ-terms (\scm{ƛ}, \scm{@}), the goals \scm{∀°} and \scm{⊃°} with
        \scm{define-dynamic-relation}, and $\Llam$ unification with universe levels and delayed
        non-patterns. Through the hooks it also makes \micro's own \scm{≠°}, \scm{absent°},
        \scm{symbol°} and \scm{number°} λ-aware; it exports no second version of them.
\end{itemize}
\noindent The design was not chosen on paper. Three prototypes were built: one that leaves \micro{}
untouched, one with a hook protocol, one with λ-terms native in the core. They were compared by three
reviews, on engineering, on fidelity to the book and on pedagogy. The hook design, ranked first by two of
the three, became the base, and about twenty grafts proposed from the other two were each adopted or
explicitly declined. Three adversarial review rounds
then attacked the result and sent thirteen, eight and twelve reports. All were real. Every defect but
one was fixed and pinned by a regression case\sidenote{A cost reported by the third round among them.}; the
exception, the cost of the scope check on first-order search once a \scm{∀°} has run, reported by the
second round, is measured and documented as a limitation. The final suite has
83 cases with 320 hand-written assertions, plus 6 cases that pin the known limits, and \micro's own 125
cases pass unchanged, with and without the module loaded: 339 cases, no failures. The protocol costs about
1\% on first-order search when no extension is loaded. All of this lives in
\texttt{doc/hol-kanren/prototype/}, and after merging it would be \texttt{src/aux.kanren.lambda.scm} and
\texttt{src/test/microkanren-lambda*.scm}.

% ==================================================================================================
\section{The whole story in brief}\label{sec:intro-brief}

This section follows one program from the book through the whole document: type inference for the untyped
λ-calculus \book{190--191}. It is two clauses of λProlog, and it uses every feature at once.

\subsection{The program, twice}

Object-level λ-terms are built with two constructors, \lp{app} for application and \lp{abs} for
abstraction. \lp{abs} takes a \emph{meta-level} abstraction: the identity is \lp{abs x\ x}. Types are
built with \lp{arr}. The relation \lp{typeof M T} says that term \lp{M} has type \lp{T}:
\begin{lprolog}
typeof (app M N) A :- typeof M (arr B A), typeof N B.
typeof (abs M) (arr A B) :- pi x\ typeof x A => typeof (M x) B.
\end{lprolog}
The second clause reads: \enquote{to type \lp{abs M}, take a new constant \lp{x}, assume that \lp{x} has
type \lp{A}, and type the body \lp{(M x)} with \lp{B}}. There is no clause for variables and no context
argument. The hypothesis \lp{typeof x A} \emph{is} the context, and it disappears when the proof of the
body is over. In the proposed module the same program reads:
\begin{scheme}
(define-dynamic-relation (typeof° m t)
  (cond°
    ((fresh° (M N B) (=° m `(app ,M ,N)) (typeof° M `(arr ,B ,t)) (typeof° N B)))
    ((fresh° (R A B) (=° m `(abs ,R)) (=° t `(arr ,A ,B))
       (∀° (x) (⊃° ((typeof° x A)) (typeof° (@ R x) B)))))))
\end{scheme}
\noindent The dictionary is short. \lp{pi x\ G} is \scm{(∀° (x) G)}, \lp{D => G} is \scm{(⊃° (D) G)}, the
application \lp{(M x)} of a logic variable is \scm{(@ R x)}, and \lp{x\ t} is \scm{(ƛ (x) t)}. The relation
is made by \scm{define-dynamic-relation} instead of \scm{define-relation} because it receives hypotheses.
The first-order structure (\scm{abs}, \scm{app}, \scm{arr}) stays ordinary Scheme data, as everywhere in
\micro.

\subsection{Running it forwards}

The book types the identity and fails on self-application \book{191}:
\begin{lprolog}
?- typeof (abs x\x) Ty.          Ty = arr T1 T1
?- typeof (abs x\ app x x) Ty.   no
\end{lprolog}
The module gives the same answers. \scm{°->list/datum} lists the answers as plain data, and \scm{α} is a
variable that the answer leaves unbound:
\begin{scheme}
(°->list/datum (fresh° (q) (typeof° `(abs ,(ƛ (x) x)) q)))
(°->list/datum (fresh° (q) (typeof° `(abs ,(ƛ (x) `(app ,x ,x))) q)))
\end{scheme}
\begin{answers}
((arr α α))
()
\end{answers}

Two rules of the book's proof search are at work, GENERIC for $\forall$ and AUGMENT for $\imp$
(\cref{sec:hh-rules}). In the sequent $\seq{\Sig}{\Prog}{G}$, $\Sig$ is the signature, the constants in
scope, and $\Prog$ is the program:
\begin{mathpar}
\inferrule*[right=\rulename{generic}]{\seq{\Sig\cup\{c\}}{\Prog}{\subst{G}{c}{x}}}{\seq{\Sig}{\Prog}{\forall x\,G}}
\quad (c\notin\Sig)
\and
\inferrule*[right=\rulename{augment}]{\seq{\Sig}{\Prog\cup\{D\}}{G}}{\seq{\Sig}{\Prog}{D\imp G}}
\end{mathpar}
Read bottom-up, GENERIC proves a universal goal by inventing a constant, and AUGMENT proves an implication
by assuming its antecedent. Both changes are local: they hold for the proof of the premise only.

\begin{tutorial}[What happens to the state while \texttt{typeof°} types λx.x]
Write $d$ for the current \emph{depth}, the number of enclosing \scm{∀°}, and $\Gamma$ for the assumed
clauses. Both live in the search state, in the slot that the protocol adds; the substitution is
\micro's, as before. Each row of the table below is one step towards the first answer, and two of them
carry the ideas. At step~5, $R$ is already bound to $\abs{x}{x}$, so \texttt{(@ R c)} reduces by β to $c$
itself: object-level substitution is meta-level β. At step~6 the dynamic relation tries its hypotheses
first, most recent first, and then its two static clauses, which both fail on $c$, since an
eigenvariable is neither an \texttt{app} nor an \texttt{abs}. The hypothesis \texttt{typeof° c A}
mentions the variable $A$ of the enclosing clause, and $A$ is \emph{shared}, not renamed. That is how
typing the body tells the outer goal what $A$ is (\cref{sec:hh-link}).
\begin{center}\small
\begin{tabular}{@{}r>{\raggedright\arraybackslash}p{5.6cm}cl@{}}
\toprule
& goal & $d$ & effect \\
\midrule
1 & \texttt{typeof° (abs (ƛ (x) x)) q} & 0 & the \texttt{app} clause fails; the \texttt{abs} clause applies\\
2 & \texttt{(abs R) = (abs (ƛ (x) x))} & 0 & first-order: $R := \abs{x}{x}$, \texttt{q := (arr A B)}\\
3 & \texttt{∀° (x)} & 1 & a new eigenvariable $c$, of level 1\\
4 & \texttt{⊃° ((typeof° c A))} & 1 & $\Gamma = \{\texttt{typeof° c A}\}$\\
5 & \texttt{typeof° (@ R c) B} & 1 & \texttt{(@ R c)} head-normalizes by β to $c$\\
6 & backchain on the hypothesis & 1 & \texttt{typeof° c A = typeof° c B}: $A$, $B$ unified\\
7 & answer leaves \texttt{⊃°} and \texttt{∀°} & 0 & $\Gamma=\emptyset$; \texttt{q = (arr A A)}, printed \texttt{(arr α α)}\\
\bottomrule
\end{tabular}
\end{center}
Self-application fails at the same step. For \texttt{(abs (ƛ (x) (app x x)))}, the body needs
$\texttt{typeof° c (arr B$'$ B)}$ and $\texttt{typeof° c B$'$}$, so the hypothesis forces
$A = \texttt{(arr B$'$ B)}$ and $A = B'$: first-order unification fails on its occurs check, as in the book.
\end{tutorial}
\noindent The book explains this program in \cref{sec:lambdatree-typeof}; \cref{sec:module-typeof} traces
it again at the level of the module's state, and \cref{sec:examples-typeof} runs it on the book's whole
session, the combinator S included.

\subsection{Running it backwards}

Ask instead for terms of type $i\to i$. λProlog's depth-first search loops on this query, because it tries
the \lp{app} clause first and descends forever. The interleaving search of \micro{} does not:
\begin{scheme}
(°->list/datum (take° 3 (fresh° (q) (typeof° q '(arr i i)))))
\end{scheme}
\begin{answers}
((abs (ƛ (x) x))
 (abs (ƛ (x) (app (abs (ƛ (x1) x1)) x)))
 (abs (ƛ (x) (app (abs (ƛ (x1) x)) x))))
\end{answers}
\noindent The second answer is $\lambda x.\,(\lambda x_1.x_1)\,x$. The printer renames the inner binder
to \texttt{x1}, apart from the outer one. The first answer is the instructive one. Now $R$ is unknown when
step~5 is reached, so backchaining on the hypothesis must solve the equation $R\,c \doteq c$, where $R$ was
created outside the \scm{∀°} and $c$ inside it. That is a \emph{higher-order} unification problem, and it
is one of the easy ones.

\begin{definition}[Levels and patterns]\label{def:intro-levels}
Every logic variable $X$ has a \emph{level} $\ell(X)$, the depth at which it was created. Every
eigenvariable $c$ has the level $\ell(c)$ of the \scm{∀°} that made it, one more than the depth outside.
\begin{enumerate}[nosep]
  \item \emph{Scope check.} A binding $X := t$, made by any \scm{=°}, is accepted only if every
        eigenvariable $c$ in $t$ has $\ell(c)\le\ell(X)$. Each unbound variable $Y$ of $t$ with
        $\ell(Y)>\ell(X)$ is then \emph{lowered}: replaced by a new variable of level $\ell(X)$ or, when
        $Y$ occurs applied, \emph{pruned} of the arguments that $X$ may not see.
  \item \emph{Pattern.} An application $F\,a_1\dots a_n$ of an unbound variable $F$ is a \emph{pattern}
        when its arguments, after head normalization, are distinct eigenvariables, each with
        $\ell(a_i)>\ell(F)$.
\end{enumerate}
The full rule, including the cases that must wait for a later binding, is in \cref{sec:module-unif}.
\end{definition}
\noindent The first condition is the \emph{scope check}. It is the operational content of the order of
quantifiers in $\exists R\,\forall c$: $R$ was chosen before $c$ existed, so its value cannot mention $c$
(\cref{sec:hh-scope,sec:lambda-unifprob}). Levels replace the book's explicit quantifier prefix, as Teyjus
does~\cite{NadathurJayaramanKwon1995}. The second condition defines Miller's fragment $\Llam$, in which
unification is decidable and has most general unifiers~\cite{Miller1991} (\cref{sec:hounif-llambda}).

\begin{example}[Scope and patterns]\label{ex:intro-scope}
Here $\ell(R)=0$ and $\ell(c)=1$. The equation $R\doteq c$ violates the scope check and fails. The equation
$R\,c\doteq c$ is a pattern: abstracting $c$ out of the right-hand side gives the most general unifier
$R := \abs{x}{x}$. In the module, with \scm{q} in the role of $R$:
\begin{scheme}
(°->list #f (fresh° (q) (∀° (x) (=° q x))))
(°->list/datum (fresh° (q) (∀° (x) (=° (@ q x) x))))
\end{scheme}
\begin{answers}
()
((ƛ (x) x))
\end{answers}
The first is the book's \lp{sigma x\ pi y\ x = y}, which is not provable \book{85}: no single term
equals every new constant.
\end{example}

The other two answers of the backwards run come from the \lp{app} clause under the binder: the body
\scm{(@ R c)} must unify with \scm{(app M N)}, where $M$ and $N$ were created inside the \scm{∀°}, at level~1.
Variable elimination (\cref{sec:hounif-algo}) builds $R := \abs{y}{\texttt{(app (M$'$ y) (N$'$ y))}}$ and
\emph{lowers} $M$ and $N$ to new variables $M'$, $N'$ of level~0, applied to $c$: $M := M'\,c$,
$N := N'\,c$. Scheme code built none of these λ-terms; unification did.

\paragraph{When an equation is not a pattern.} Some equations have no most general unifier. The book's
$\forall x.\,F\,x\,x\doteq x$ \book{208} has two incomparable solutions, $\abs{a}{\abs{b}{a}}$ and
$\abs{a}{\abs{b}{b}}$, and $F\,x\,x$ is not a pattern, since its arguments repeat. The module does what
Teyjus does, \enquote{dynamic $\Llam$} (\cref{sec:hounif-dynamic}): it suspends the equation, retries
it after every later binding, which may have turned it into a pattern, and reports it in the answer if it
is still there at the end. Under \scm{°->pretty}, which prints constraints in the style of αKanren, the residual equation is
closed over $x$ when it leaves the scope of the \scm{∀°}. A later binding decides it:
\begin{scheme}
(°->pretty (fresh° (q) (∀° (x) (=° (@ q x x) x))))
(°->list/datum (fresh° (q) (∀° (x) (=° (@ q x x) x)) (=° q (ƛ (a b) b))))
\end{scheme}
\begin{answers}
((α : (= (ƛ (x) (@ α x x)) (ƛ (x) x))))
((ƛ (a b) b))
\end{answers}

\subsection{What the engine had to learn}

None of the above runs on today's \micro. \Cref{tab:intro-concepts} lists what the example needed, where
the book explains it, and where this document builds it.
\begin{table*}[h]
\small
\begin{tabularx}{\linewidth}{@{}>{\raggedright\arraybackslash}p{4.4cm}>{\raggedright\arraybackslash}X>{\raggedright\arraybackslash}p{4.3cm}@{}}
\toprule
concept & in the example & theory \,/\, in the module\\
\midrule
λ-tree syntax, α-equivalence & \texttt{(abs (ƛ (x) x))}, binders printed apart & \cref{ch:lambdatree} \,/\, \cref{ch:module}\\
β-reduction, application & \texttt{(@ R x)} is \texttt{x} once \texttt{R := (ƛ (x) x)} & \cref{ch:lambda} \,/\, \cref{ch:module}\\
GENERIC, eigenvariables & \texttt{(∀° (x) …)} & \cref{ch:hh} \,/\, \cref{ch:design,ch:module}\\
AUGMENT, dynamic relations & \texttt{(⊃° ((typeof° x A)) …)} & \cref{ch:hh} \,/\, \cref{ch:module}\\
scope check and levels & \texttt{R = x} fails & \cref{ch:hh,ch:hounif} \,/\, \cref{ch:module}\\
pattern unification & \texttt{R x = x} gives \texttt{R := λx.x} & \cref{ch:hounif} \,/\, \cref{ch:module}\\
dynamic $\Llam$, residuals & \texttt{F x x = x} waits & \cref{ch:hounif} \,/\, \cref{ch:module}\\
interleaving search & typing backwards terminates & \cref{ch:horn} \,/\, \cref{ch:micro}\\
\bottomrule
\end{tabularx}
\caption{The running example, concept by concept.}\label{tab:intro-concepts}
\end{table*}

The state of \micro{} has a fixed set of fields, and its unifier, walks and reifier are closed
(\cref{sec:micro-ext}). The proposal opens them once, generically, and puts everything λ-specific in a
client module. In the state's new slot the module keeps the depth, the level of each variable, the
assumed clauses $\Gamma$, the suspended constraints (equations, and types or absentos on flexible
terms), and the disequalities that involve λ-terms or eigenvariables. Through the hooks it receives every
unification problem that contains one of its terms, it checks the scope of the bindings that \scm{=°}
makes, and it takes part in printing. Because the hooks sit inside \micro's own \scm{=°}, every relation
built on \scm{=°}, such as \micro's \scm{cons°} or the user's own, respects the scope of \scm{∀°} with no
change. A prototype that wrapped \micro{} from the outside instead let eigenvariables escape through exactly
such relations (\cref{ch:design}).

\begin{kanren}[Why not just use Scheme closures?]
A Scheme procedure can stand for a binder as long as it is only \emph{called}: \lp{mappred} and continuation
passing work on today's \micro{} (\cref{sec:hoprog-free}). The backwards run above needs more: unification
must look \emph{inside} the abstraction, compare it modulo α, and build one. A closure can do none of these,
so \scm{ƛ} builds a data structure, a locally nameless term whose binders are de Bruijn
indices~\cite{deBruijn1972,Chargueraud2012}, which prints with the user's names.
\end{kanren}

\begin{pitfall}[Hypotheses need a dynamic relation]
In λProlog every predicate can receive clauses through \lp{=>}. In the module only a relation made by
\scm{define-dynamic-relation} can, because a relation made by \scm{define-relation} is a closed Scheme
procedure. Assuming a clause for it is an error, not a silent no-op:
\begin{scheme}
(define-relation (static° x) (=° x 1))
(define-dynamic-relation (dynamic° x) (=° x 1))
(°->list/datum (fresh° (q) (⊃° ((static° 2)) (static° q))))   ; an error
(°->list/datum (fresh° (q) (⊃° ((dynamic° 2)) (dynamic° q))))
\end{scheme}
\begin{answers}
Error: (⊃°) not a dynamic relation, see define-dynamic-relation: #<procedure (static° x)>
(2 1)
\end{answers}
\noindent The first line is what \texttt{csi} prints for the error, which stops the script; the last
query was run on its own. In the second answer the hypothesis comes first, then the
static clause, as in the book \book{78}.
\end{pitfall}

% ==================================================================================================
\section{Map of the document}\label{sec:intro-map}

The document has three parts and two appendices. Part~I follows the book, Part~II the literature, and
Part~III builds the proposal.

\subsection*{Part I: the book}

\paragraph{\Cref{ch:horn}, Logic programming as proof search.} The book's Introduction and Chapters~1--2:
typed first-order terms, first-order unification, the search semantics of goals, and backchaining on Horn
clauses (\fohc). It establishes that pure miniKanren \emph{is} \fohc{} without types and with an interleaving
search, and it shows where the interleaving finds answers that depth-first search misses.

\paragraph{\Cref{ch:hh}, Hereditary Harrop formulas.} Chapter~3 of the book: AUGMENT and GENERIC, hypothetical
reasoning on a course database, scoped constants, why $\exists\forall$ and $\forall\exists$
differ, and how assumed clauses share variables with their context. The logic is intuitionistic, indeed
minimal, not classical. Its last section lists what $\forall$ and $\imp$ ask of a miniKanren: the program
context must live in the search state, and scoped constants need a scope check.

\paragraph{\Cref{ch:lambda}, Typed λ-terms and unification problems.} Chapter~4: simply typed λ-terms,
conversion by α, β and η as the equality of the logic, normal forms, unification problems as quantified formulas,
raising, and why higher-order unification is undecidable (Post's correspondence problem, Hilbert's tenth
problem).

\paragraph{\Cref{ch:hoprog}, Higher-order programming.} {\emergencystretch=2em Chapters~5--6: rigid and flexible atoms, the
languages \hohc{} and \hohh{}, predicate variables (\lp{mappred}, folds, continuation passing),
function variables solved by unification (functional difference lists), and modules as logic. It sorts
the book's programs into those that run on today's \micro{} (they only \emph{call} a predicate) and those
that need λ-terms (they look \emph{inside} one).\par}

\paragraph{\Cref{ch:lambdatree}, λ-tree syntax.} {\emergencystretch=2em Chapter~7, the chapter that motivates the proposal:
object-level binders as meta-level abstractions, substitution as β, the \emph{mobility} of binders
(term-level λ, then \lp{pi}, then an eigenvariable), and the book's programs over λ-terms and formulas:
evaluators, type inference, de~Bruijn translation, normal forms, prenex forms, copy clauses, and
substitution run backwards. It ends with the $\Llam$ subset.\par}

\paragraph{\Cref{ch:hounif}, Unification of λ-terms and higher-order patterns.} Chapter~8: Huet's procedure
and its matching trees, the pattern fragment, the pattern unification algorithm as one walk of variable
elimination (inversion, pruning, occurs and scope checks), traces of the book's examples, Nipkow's
de~Bruijn algorithm with a running Scheme transliteration, and the pragmatics of real systems: universe
levels in place of raising, and delayed non-patterns.

\paragraph{\Cref{ch:apps}, Applications.} Chapters~9--11: provers for intuitionistic and classical logic,
natural deduction with proof terms, tactics, typing and evaluation of functional programs, partial
evaluation, CPS conversion, and the π-calculus. For each application it records which of $\forall$, $\imp$
and λ-term unification it needs, and which parts run on \micro{} today.

\subsection*{Part II: the literature}

\paragraph{\Cref{ch:lkanren}, Higher-order miniKanren.} Ma and Friedman 2021 in full: grafting versus
substitution, binder merging and its \enquote{true = false} hazard, static closures, and a critical assessment
with three counter-examples. Then λKanren 2020, nominal unification and αKanren, eigenvariables and negation
inside miniKanren, and the implementation techniques of Miller, Nipkow, Dowek et al., Nadathur et al.,
Teyjus, ELPI and Makam. It ends with a comparison of the design options (\cref{sec:lkanren-compare}) and
deliberately does not choose.

\subsection*{Part III: the proposal}

\paragraph{\Cref{ch:micro}, Anatomy of \micro.} {\emergencystretch=2em The engine as it is today, at commit \texttt{f53bdcb}: terms
and variables, the six-field state, the walk, unification, the stores D (disequalities), A (type tags) and T
(absento), the interleaving search, reification and aggregation. It closes with what is open to extension and
what is closed, and with three tricks that need no edit (including working AUGMENT and GENERIC prototypes)
together with their limits.\par}

\paragraph{\Cref{ch:design}: requirements, design space, protocol.} The nine acceptance items A1--A9
collected from Parts~I and~II (the scope of $\forall$ at every binding, AUGMENT with the book's clause
order, type inference, λ-tree syntax, $\Llam$ with delay, λProlog's answers where Ma and Friedman differ,
the book's programs, \micro's stores on λ-terms, and \micro's own suites unchanged). Then the design space,
the four meanings of \scm{=°} on λ-terms left open by \cref{sec:lkanren-compare}, of which pattern
unification with levels is retained; the three prototypes (a layer on top of \micro, hooks in \micro, a
native engine) and the three reviews that ranked them; and the extension protocol that came out of it:
the state's seventh field, the registry, the thirteen hooks with their contracts and call sites, the
invariants every client must keep, and its cost. It ends with the history of the adversarial review rounds.

\paragraph{\Cref{ch:module}: the module.} \texttt{(aux kanren lambda)} in detail: locally nameless terms
built by \scm{ƛ} and \scm{@}; the extension's state (depth, levels, $\Gamma$, suspended constraints,
λ-disequalities), copied and never mutated; the unification rules with variable elimination, the eager
scope check, the fast paths and the wake-up loop, with worked traces; \scm{∀°} and \scm{⊃°}, including
the discharge of constraints at the exit of a scope and backchaining in \scm{define-dynamic-relation};
disequality, absento and types on λ-terms; reification; and the API reference.

\paragraph{\Cref{ch:examples}: the book's examples in the module.} The core programs of Part~I, each beside its
λProlog original and with its actual output: α, β and shadowing; the scope of eigenvariables; the course
database, jars and bugs, and reverse with assumed clauses; type inference; call-by-name evaluation, copy and
substitution, existential generalization; de~Bruijn translation and prenex forms; \lp{mappred} and
palindromes with functional difference lists; Fischer's CPS transformation; and the unification examples of
the book's Chapter~8 one by one. It also runs Ma and Friedman's queries and explains each difference, and
closes with a table of all the examples.

\paragraph{\Cref{ch:roadmap}: tests, limitations, performance, merge plan.} {\emergencystretch=2em How the proposal is tested
(three layers of tests, and \micro's suites run twice); nine documented limitations and deviations from
λProlog, each pinned by a test (among them: only dynamic relations take hypotheses, terms are untyped,
non-patterns are suspended and never enumerated, answers come in \micro's interleaving order); the
performance measurements; a merge in three stages (first the protocol, then the module and its suites,
then the show extension); six later work items; and nine questions only the maintainer of \aux{} can
answer.\par}

\subsection*{Appendices}

\paragraph{\Cref{app:rules}, Reference card.} The formal rules on a few pages: λProlog syntax and its
logical reading, the goal and clause grammars, uniform proofs, typing, λ-conversion, raising, first-order
unification, Huet's procedure, pattern unification and universe levels, each with its page in the book.

\paragraph{\Cref{app:glossary}, Glossary.} The vocabulary of the three communities (proof theory,
unification, miniKanren) in one alphabetical list, with pointers to the chapters and to the book, and a
bridge between their names for the same idea.

% ==================================================================================================
\section{Reading paths}\label{sec:intro-paths}

The chapters are written to be read in order, but they need not be. \Cref{tab:intro-paths} summarizes the
three paths below; page counts are approximate.

\paragraph{The fast path.} The minimum to understand the proposal, about 45 pages before
\cref{ch:design,ch:module}:
\begin{enumerate}[nosep]
  \item \cref{sec:intro-brief} of this chapter (4 pages);
  \item the two rules, the scope condition, and how quantification links goals and clauses:
        \cref{sec:hh-rules,sec:hh-scope,sec:hh-link} (5 pages);
  \item λ-tree syntax:
        \cref{sec:lambdatree-encoding,sec:lambdatree-beta,sec:lambdatree-mobility,sec:lambdatree-typeof}
        (9 pages);
  \item what rigid and flexible terms and the normal form of a unification problem are:
        \cref{sec:hounif-intro,sec:hounif-problems} (the introduction says enough about imitation and
        projection to skip \cref{sec:hounif-huet}); then patterns, their unification, the worked traces
        and what real systems do: \cref{sec:hounif-llambda,sec:hounif-algo,sec:hounif-traces,sec:hounif-pragmatics}
        (13 pages in all);
  \item the options: \cref{sec:lkanren-compare,sec:lkanren-summary} (3 pages; the summary defines each
        term the comparison uses);
  \item the engine's state, its constraint stores and its extension points:
        \cref{sec:micro-state,sec:micro-stores,sec:micro-ext} (8 pages; skip \cref{sec:micro-stores} if you
        know micro's D/A/T chain);
  \item then \cref{ch:design,ch:module} (about 60 pages), the limitations in \cref{sec:roadmap-limits},
        and the staged merge in \cref{sec:roadmap-merge} (2 pages), which says how to build and load
        the edited micro.
\end{enumerate}

\paragraph{The λProlog path.} Part~I alone, \cref{ch:horn} to \cref{ch:apps}, about 185 pages, is a
self-contained reading of the book. The \emph{Kanren angle} boxes can be skipped. \Cref{app:rules} serves as
a reference card along the way.

\paragraph{The implementer path.} For whoever will merge or extend the code:
\cref{ch:hounif} (the algorithm), \cref{ch:micro} (the engine), \cref{ch:design} (the protocol),
\cref{ch:module} (the module), and \cref{ch:roadmap} (tests, costs, merge plan), about 120 pages, with
\cref{sec:rules-pattern,sec:rules-levels} at hand. \Cref{ch:examples} then serves as a set of worked tests.

\begin{table}[htb]
\centering\small
\caption{Three reading paths. A filled bullet means read in full, a hollow one read the sections named in
the text. Pages are those of this document.}\label{tab:intro-paths}
\begin{tabular}{@{}lrccc@{}}
\toprule
chapter & pages & fast & λProlog & implementer\\
\midrule
\cref{ch:intro}, Introduction & 12 & $\circ$ & $\bullet$ & $\bullet$\\
\cref{ch:horn}, proof search & 28 & & $\bullet$ & \\
\cref{ch:hh}, hereditary Harrop & 24 & $\circ$ & $\bullet$ & \\
\cref{ch:lambda}, λ-terms & 22 & & $\bullet$ & \\
\cref{ch:hoprog}, higher-order programming & 30 & & $\bullet$ & \\
\cref{ch:lambdatree}, λ-tree syntax & 30 & $\circ$ & $\bullet$ & \\
\cref{ch:hounif}, pattern unification & 24 & $\circ$ & $\bullet$ & $\bullet$\\
\cref{ch:apps}, applications & 28 & & $\bullet$ & \\
\cref{ch:lkanren}, the literature & 30 & $\circ$ & & \\
\cref{ch:micro}, the engine & 22 & $\circ$ & & $\bullet$\\
\cref{ch:design}, design and protocol & 26 & $\bullet$ & & $\bullet$\\
\cref{ch:module}, the module & 32 & $\bullet$ & & $\bullet$\\
\cref{ch:examples}, the examples & 28 & & & $\bullet$\\
\cref{ch:roadmap}, status and roadmap & 18 & $\circ$ & & $\bullet$\\
\cref{app:rules}, reference card & 12 & & $\circ$ & $\circ$\\
\bottomrule
\end{tabular}
\end{table}


% ==================================================================================================
\section{Conventions}\label{sec:intro-conventions}

\paragraph{Boxes.} Three kinds of box interrupt the text, and each can be skipped on a first reading.
A \textbf{Tutorial} box walks through a computation step by step: a trace of unification, of a $\forall$ or
$\imp$ step, of the changes of a state (see the box in \cref{sec:intro-brief}). A \textbf{Kanren angle} box
says how a concept maps to \micro, what runs there today, and what the proposal changes. A \textbf{Pitfall}
box marks a trap: one the book stresses, or one that a Scheme programmer is likely to fall into.

\paragraph{Code.} Three environments, told apart by the color of the rule on their left:
\begin{lprolog}
typeof (abs M) (arr A B) :- pi x\ typeof x A => typeof (M x) B.   % λProlog
\end{lprolog}
\begin{scheme}
(°->list/datum (fresh° (q) (typeof° `(abs ,(ƛ (x) x)) q)))       ; aux Scheme
\end{scheme}
\begin{answers}
((arr α α))
\end{answers}
\noindent The last one, \texttt{answers}, is always the \emph{actual output} of a run of the Scheme code
above it, pasted verbatim, never written by hand. Inline, \lp{pi x\ G} is λProlog and \scm{(∀° (x) g)} is
Scheme. Answers come in three forms. \scm{°->list/datum} gives plain data, comparable with \scm{equal?}.
\scm{°->list #f} gives \micro's own form, a closed λ-expression whose parameters are the answer's unbound
variables, such as \scm{(λ () (ƛ (x) x))}. \scm{°->pretty}, from \texttt{(aux kanren micro show)}, gives
the αKanren-style view with constraints, such as \scm{(α : (= …))}. Reified variables print as \scm{α}, \scm{β}, \scm{γ}; a working variable that a
constraint keeps prints as \scm{_3}. Names follow \aux: goals end in \scm{°}, their applicative forms in
\scm{ª}, term constructors carry no suffix (\scm{ƛ}, \scm{@}), and internals carry the prefix
\scm{μkanren-}.

\paragraph{Pages of the book.} \book{n} refers to the \emph{printed} page $n$ of~\cite{MillerNadathur2012}, the
number at the top of the page, which is the PDF page minus 16.
\begin{pitfall}[Two page numberings]
The comments in the prototype's files (\texttt{examples.scm}, the test suite, the module's header) cite the
book as \texttt{p.N} with the \emph{PDF} page. So \texttt{p.207} in \texttt{examples.scm}, next to the
\texttt{typeof} queries, is \book{191} here.
\end{pitfall}

\paragraph{Notation.} A sequent $\seq{\Sig}{\Prog}{G}$ asks for a proof of goal $G$ from program $\Prog$ over
signature $\Sig$. Implication is written $\imp$, as in the book, and $D\imp G$ is \lp{D => G}. The logics are
\fohc{} and \fohh{} (first-order Horn and hereditary Harrop) and \hohc{} and \hohh{} (their higher-order
versions). $\Llam$ is the pattern fragment, $\eqb$ is λ-convertibility, $\subst{t}{s}{x}$ substitutes $s$ for $x$
in $t$, and $\doteq$ is an equation to be solved. $\ell(X)$ and $\ell(c)$ are the levels of
\cref{def:intro-levels}. Logic variables are upper case ($F$, $X$, $R$), eigenvariables lower case ($c$,
$x$, $y$). Rules of proof search are set in small capitals (\rulename{generic}, \rulename{augment},
\rulename{backchain}).

\paragraph{Where the outputs come from.} {\emergencystretch=2em Snippets in Parts~I and~II, and in \cref{ch:micro}, run on today's
\micro, at commit \texttt{f53bdcb} of the \texttt{aux-\allowbreak lambdakanren} branch, unless a box says otherwise; some
define a small throwaway interpreter first (\cref{sec:hh-interp}). The outputs of the proposal, in this
chapter and from \cref{ch:design} on, come from the prototype as described in Part~III: \micro{} with the
protocol applied (\texttt{micro.diff}), the show module with \texttt{show.diff}, and
\texttt{(aux kanren lambda)}, all in \texttt{doc/\allowbreak hol-kanren/\allowbreak prototype/}. Every snippet was run by
\texttt{csi} from CHICKEN~6.0.1pre1 (rev \texttt{55615603}), as installed; the
benchmarks were compiled with \mbox{\texttt{csc -O2 -d1}}.\marginnote{The document was typeset with XeLaTeX, through Tectonic
0.17.0, in Libertinus and DejaVu Sans Mono.}\par}

% ==================================================================================================
\section{Summary}\label{sec:intro-summary}
\begin{itemize}
  \item The goal is to bring λProlog's universal and hypothetical goals, λ-tree syntax and higher-order pattern
        unification into \micro, following Miller and Nadathur's book (Part~I) and the miniKanren literature,
        Ma and Friedman 2021 first (Part~II).
  \item The outcome (Part~III) is a generic extension protocol in \micro{} plus the client module
        \texttt{(aux kanren lambda)}. Three prototypes were compared by three reviews, and three adversarial
        rounds hardened the result: 89 λ cases and \micro's 125 cases, run with and without the module,
        pass, and the protocol costs about 1\% on first-order search.
  \item The running example is \lp{typeof}. \scm{∀°} gives the body an eigenvariable one level deeper, and
        \scm{⊃°} assumes its type. Backchaining on that hypothesis unifies the types, and the assumption
        disappears with the answer.
  \item Run backwards, the same program needs $R\,c\doteq c$: a pattern, solved by $R:=\abs{x}{x}$, while
        $R\doteq c$ fails the scope check. Non-patterns wait and are reported as residual equations.
  \item Three reading paths: a fast path of about 45 pages before the design and the module, the λProlog
        path (Part~I), and the implementer path.
  \item Answers are actual outputs; \book{n} is a printed page; the prototype's comments use PDF pages
        (printed $+16$).
\end{itemize}
